\documentclass[11pt]{article}

\usepackage[a4paper,margin=2.5cm]{geometry}
\usepackage[ngerman]{babel}
\usepackage{amsmath,amssymb,amsthm,amsfonts}
\usepackage{mathtools}
\usepackage{hyperref}
\usepackage{enumitem}

\title{Priml\"ucken als Rotationen auf einem zyklischen Vier-Zustands-Raum\\
\large (Revision~V2: Br\"ucke zu E-ENT-003)}
\author{Thomas Hoffbauer}
\date{}

\newtheorem{theorem}{Satz}
\newtheorem{lemma}{Lemma}
\newtheorem{corollary}{Korollar}
\newtheorem{proposition}{Proposition}

\theoremstyle{definition}
\newtheorem{definition}{Definition}
\newtheorem{remark}{Bemerkung}

\begin{document}

\maketitle

\begin{abstract}
Wir betrachten Priml\"ucken $g_n=p_{n+1}-p_n$ zwischen aufeinanderfolgenden Primzahlen $>3$ modulo $12$. Die zul\"assigen Restklassen sind $\{1,5,7,11\}=(\mathbb{Z}/12\mathbb{Z})^\times$ mit $(\mathbb{Z}/12\mathbb{Z})^\times\cong C_2\times C_2$; verwendet wird jedoch nur eine zus\"atzlich gew\"ahlte zyklische Ordnung dieser vier Klassen. Der Hauptsatz zeigt, dass das zugeh\"orige Rotationsinkrement auf einem gerichteten Kreisgraphen $C_4$ ausschlie{\ss}lich von $g_n\bmod 12$ abh\"angt. Daraus folgt eine wohldefinierte Abbildung der sechs geraden L\"uckenklassen modulo $12$ auf vier Rotationszust\"ande. Die Konstruktion liefert damit eine exakte kombinatorische Kodierung von Priml\"ucken als diskrete Rotationsdynamik.

\emph{Revision~V2.} Als optionale, strikt abgesetzte Modellschicht verweisen wir auf die TUR-Invariante
$\mathcal{Q}_{BC}\cdot S_{BC}(\phi)=\Lambda/(T\cdot\gamma)\ge\mathcal{C}_{Q_0}$
(Evidenz-ID \texttt{E-ENT-003}, Projekt \emph{Entropie}, Tag \texttt{v0.1-tur-bc}).
Diese Schicht ersetzt weder den kombinatorischen Hauptsatz noch behauptet sie eine thermodynamische Identifikation.
\end{abstract}

\section{Einleitung}

Priml\"ucken werden traditionell als arithmetische Gr\"o{\ss}en untersucht,
\[
g_n = p_{n+1}-p_n,
\]
wobei $p_n$ und $p_{n+1}$ aufeinanderfolgende Primzahlen sind.

Das Ziel dieser Notiz ist zu zeigen, dass dieselbe Information eine nat\"urliche geometrische Interpretation besitzt.
Anstatt Priml\"ucken als Abst\"ande auf der Zahlengeraden zu betrachten, interpretieren wir sie als Erzeuger von Rotationen auf einem zyklischen Vier-Zustands-System, das aus den zul\"assigen Restklassen modulo zw\"olf gewonnen wird.

Der mathematische Anspruch dieser Darstellung liegt in einer sauberen, exakten Umcodierung, nicht in einer neuen tiefen Struktur der Primzahlen. Insbesondere wird keine neue Gruppenstruktur eingef\"uhrt; vielmehr wird auf einer endlichen Vierermenge eine zus\"atzliche zyklische Ordnung gew\"ahlt, die die L\"uckenstatistik in eine dynamische Sprache \"ubersetzt.

\section{Der Restklassen-Zustandsraum}

Sei
\[
U_{12}
=
(\mathbb{Z}/12\mathbb{Z})^\times
=
\{1,5,7,11\}.
\]
Jede Primzahl $p>3$ geh\"ort genau zu einer dieser Restklassen.
Als Gruppe gilt $U_{12}\cong C_2\times C_2$, also ist $U_{12}$ nicht zyklisch.
Die Verwendung eines Zyklus bezieht sich hier nur auf eine zus\"atzliche Ordnung dieser vier Elemente.

\begin{definition}
Definiere die Zustandsabbildung
\[
\rho:U_{12}\rightarrow \mathbb{Z}_4
\]
durch
\[
\rho(1)=0,\qquad
\rho(5)=1,\qquad
\rho(7)=2,\qquad
\rho(11)=3.
\]
\end{definition}

Damit bestimmt jede Primzahl $p>3$ einen Zustand
\[
s(p)=\rho(p\bmod 12)\in\mathbb{Z}_4.
\]
Geometrisch identifiziert dies die zul\"assigen Restklassen mit dem gerichteten Kreisgraphen
\[
C_4:\qquad 0\rightarrow1\rightarrow2\rightarrow3\rightarrow0.
\]
\"Aquivalent dazu w\"ahlen wir auf den urspr\"unglichen Restklassen die zyklische Ordnung
\[
1\to5\to7\to11\to1.
\]

Aus historischen Gr\"unden verwenden wir gelegentlich auch die Bezeichnungen
\[
0=E,\qquad 1=A,\qquad 2=B,\qquad 3=C.
\]

\section{Priml\"ucken und Rotationsinkremente}

Seien $p_n<p_{n+1}$ aufeinanderfolgende Primzahlen und
\[
g_n=p_{n+1}-p_n.
\]
Dann gilt f\"ur $p_n>3$:
\[
g_n\equiv 0,2,4,6,8,10 \pmod{12}.
\]

\begin{definition}
Das zu $g_n$ geh\"orige Rotationsinkrement ist
\[
R_n=s(p_{n+1})-s(p_n)\pmod 4.
\]
\end{definition}

\begin{proposition}
Sei $\rho:U_{12}\to\mathbb{Z}_4$ die gew\"ahlte zyklische Nummerierung und
f\"ur $d\in\{0,2,4,6,8,10\}$ die Abbildung
\[
T_d:U_{12}\to U_{12},\qquad T_d(x)\equiv x+d\pmod{12}.
\]
Dann ist jedes $T_d$ eine Permutation von $U_{12}$, und es gibt genau ein $r(d)\in\mathbb{Z}_4$ mit
\[
\rho(T_d(x))\equiv \rho(x)+r(d)\pmod 4
\quad\text{f\"ur alle }x\in U_{12}.
\]
Explizit gilt
\[
r(0)=0,\quad r(2)=r(4)=1,\quad r(6)=2,\quad r(8)=r(10)=3.
\]
\end{proposition}

\begin{proof}
Da $d$ gerade ist, bildet $x\mapsto x+d\pmod{12}$ die vier ungeraden Klassen
$U_{12}=\{1,5,7,11\}$ wieder auf sich ab; auf einer endlichen Menge ist dies eine Permutation.
Die Werte von $r(d)$ ergeben sich durch direkte Auswertung von $T_d$ auf den vier Klassen:
\[
\begin{array}{c|c}
d \bmod 12 & \text{\"Ubergangsmuster } s\mapsto s+r(d)\ (\bmod 4)\\
\hline
0   & 0\to0,\ 1\to1,\ 2\to2,\ 3\to3\\
2   & 0\to1,\ 1\to2,\ 2\to3,\ 3\to0\\
4   & 0\to1,\ 1\to2,\ 2\to3,\ 3\to0\\
6   & 0\to2,\ 1\to3,\ 2\to0,\ 3\to1\\
8   & 0\to3,\ 1\to0,\ 2\to1,\ 3\to2\\
10  & 0\to3,\ 1\to0,\ 2\to1,\ 3\to2
\end{array}
\]
Damit ist die Formel f\"ur $r(d)$ gezeigt.
\end{proof}

\section{Darstellungssatz f\"ur Rotationen}

\begin{theorem}
Seien $p_n<p_{n+1}$ aufeinanderfolgende Primzahlen gr\"o{\ss}er als drei.
Dann h\"angt das Rotationsinkrement
\[
R_n=s(p_{n+1})-s(p_n)\pmod 4
\]
nur von der Restklasse der Priml\"ucke modulo zw\"olf ab.
Genauer gilt
\[
R_n=\Phi(g_n\bmod 12),
\]
wobei
\[
\Phi:\{0,2,4,6,8,10\}\rightarrow\mathbb{Z}_4
\]
durch
\[
\Phi(0)=0,\qquad
\Phi(2)=1,\qquad
\Phi(4)=1,\qquad
\Phi(6)=2,\qquad
\Phi(8)=3,\qquad
\Phi(10)=3.
\]
\end{theorem}

\begin{proof}
Setze $d=g_n\bmod 12\in\{0,2,4,6,8,10\}$.
Nach der vorigen Proposition gilt
\[
\rho(T_d(x))\equiv \rho(x)+r(d)\pmod 4
\quad\text{f\"ur alle }x\in U_{12}.
\]
Setzt man $x=p_n\bmod 12$, so ist $T_d(x)\equiv p_n+g_n\equiv p_{n+1}\bmod 12$.
Folglich gilt
\[
s(p_{n+1})=\rho(p_{n+1}\bmod 12)=\rho(T_d(x))\equiv \rho(x)+r(d)=s(p_n)+r(d)\pmod 4,
\]
woraus $R_n=s(p_{n+1})-s(p_n)\equiv r(d)\pmod 4$ folgt.
Die Zuordnung $\Phi(d)=r(d)$ ergibt exakt die behaupteten Werte.
\end{proof}

\section{Interpretation}

\begin{remark}
Der Begriff \emph{Rotation} wird hier graphentheoretisch verwendet:
Gemeint ist die orientierte Verschiebung auf dem gerichteten Kreisgraphen $C_4$,
nicht eine geometrische Rotation im euklidischen Raum.
\end{remark}

Der Satz liefert eine exakte Faktorisierung
\[
\text{Priml\"ucken}
\longrightarrow
\mathbb{Z}/12\mathbb{Z}
\longrightarrow
C_4.
\]
Damit erzeugt jede Priml\"ucke eine eindeutige Rotation auf dem Kreisgraphen.

Die Zuordnung lautet:
\[
\begin{array}{c|c|c}
g_n\bmod 12 & R_n & \text{Interpretation}\\
\hline
0 & 0 & \text{Keine Bewegung}\\
2,4 & 1 & \text{Vorw\"artsrotation}\\
6 & 2 & \text{Halbdrehung}\\
8,10 & 3 & \text{R\"uckw\"artsrotation}
\end{array}
\]

Folglich induziert die Primzahlfolge einen diskreten Rotationsfluss
\[
(s(p_n))_{n\ge 1}
\]
auf dem Kreisgraphen $C_4$.

\section{Primzahlvierlinge}

Betrachte einen Primzahlvierling
\[
(p,p+2,p+6,p+8).
\]
Sein Restklassenmuster modulo zw\"olf ist notwendig
\[
(5,7,11,1)
\]
oder
\[
(11,1,5,7).
\]

\begin{corollary}
Jeder Primzahlvierling induziert die Rotationsfolge
\[
(1,1,1),
\]
also drei aufeinanderfolgende Vorw\"artsrotationen auf $C_4$.
\end{corollary}

\begin{proof}
Die aufeinanderfolgenden L\"ucken sind
\[
2,\;4,\;2.
\]
Jede davon geh\"ort zur Klasse
\[
\{2,4\},
\]
die $R_n=1$ entspricht.
\end{proof}

\section{Dynamische Sicht}

Die Primzahlfolge
\[
(p_n)
\]
bestimmt daher eine Trajektorie
\[
(s(p_n))
\]
auf dem endlichen Zustandsraum $C_4$.
Priml\"ucken wirken als Bewegungserzeuger, w\"ahrend Primzahlkonstellationen als koh\"arente Rotationsmuster erscheinen.

Dies ergibt eine geometrische Neuformulierung der Priml\"uckenstatistik:
\[
\boxed{
\text{Priml\"ucken modulo }12
\Longleftrightarrow
\text{Rotationen auf }C_4.
}
\]

\section{Ausblick}

Ein naheliegender n\"achster Schritt ist eine nichttriviale dynamische Aussage \"uber die induzierte Folge $(R_n)$. Als erste Leitgr\"o{\ss}e bietet sich das Verh\"altnis
\[
M_N
:=
\frac{\#\{1\le n\le N:\,R_n=1\}}
{\#\{1\le n\le N:\,R_n=3\}}.
\]
\emph{Offene Alternative.} Entweder gilt
\[
\lim_{N\to\infty} M_N = 1,
\]
oder es existiert ein Grenzwert $L\neq1$ mit
\[
\lim_{N\to\infty} M_N = L.
\]
Als empirisch vorsichtigere Zwischenhypothese kann man
\[
M_N>1
\]
auf gro{\ss}en endlichen Bereichen testen.

Dar\"uber hinaus ergeben sich die folgenden konkreten Forschungsvorschl\"age:
\begin{enumerate}[label=\textbf{F\arabic*:},leftmargin=*,itemsep=0.35em]
\item \textbf{\"Ubergangsmatrix und Markov-N\"aherung.}
Definiere f\"ur $i,j\in\mathbb{Z}_4$ die empirischen \"Ubergangszahlen
\[
N_{ij}(X):=\#\{n:\ p_n\le X,\ s(p_n)=i,\ s(p_{n+1})=j\},
\]
und die normierte Matrix $P_X=(p_{ij}(X))$ mit
\[
p_{ij}(X):=\frac{N_{ij}(X)}{\sum_{k}N_{ik}(X)}.
\]
Zu untersuchen ist die Konvergenz von $P_X$ und eine m\"ogliche Asymmetrie zwischen Vorw\"arts- und R\"uckw\"arts\"uberg\"angen.

\item \textbf{H\"aufigkeiten der Rotationsklassen.}
Setze
\[
\nu_r(X):=\#\{n:\ p_n\le X,\ R_n=r\}\qquad (r\in\mathbb{Z}_4).
\]
Die zentralen Fragen sind Grenzverh\"altnisse wie $\nu_1(X)/\nu_3(X)$, sowie Fehlerterme gegen ein Referenzmodell ohne Bias.

\item \textbf{Vergleich mit Zufallsmodellen.}
Vergleiche die beobachteten Werte von $M_N$, $\nu_r(X)$ und $P_X$ mit
geeigneten Cram\'er-artigen oder entkoppelten Restklassenmodellen.
Damit l\"asst sich trennen, welche Effekte rein lokal-modular sind und welche auf tieferer Korrelation beruhen.

\item \textbf{H\"ohere Konstellationen.}
Verallgemeinere das Vierlings-Korollar auf l\"angere Muster (z.\,B. Quintuplets,
Sextuplets, admissible constellations) und klassifiziere die zugeh\"origen
Rotationsw\"orter in $\mathbb{Z}_4$.
Dies f\"uhrt zu kombinatorischen Invarianten f\"ur Primzahlkonstellationen.

\item \textbf{Verallgemeinerung des Moduls.}
Ersetze $12$ durch ein allgemeineres Modul $m$ mit fixierten zul\"assigen Restklassen.
Die Frage ist, wann eine analoge Kodierung auf einem endlichen gerichteten Graphen
entsteht und welche dynamischen Gr\"o{\ss}en dann stabil bleiben.
\end{enumerate}

Diese Agenda macht aus der vorliegenden Kodierung ein systematisches Forschungsprogramm,
in dem numerische Evidenz, kombinatorische Struktur und analytische Heuristik unmittelbar zusammenwirken.

\section{Erweiterung: Transversale Entropie-Kopplung an den $BC$-Unterraum ($E$-ENT-$003$)}
\label{sec:e-ent-003}

Im Kontext des erweiterten Evidenzregisters wird die diskrete Rotationsdynamik auf $C_4$
\"uber das transversale Segment
$\mathcal{H}_{BC}=\mathrm{span}(B,C)=\mathrm{span}(j,k)$
an die kontinuierliche stochastische Str\"omung gekoppelt.
Der kombinatorische Hauptsatz ($R_n=\Phi(g_n\bmod 12)$) bleibt davon unber\"uhrt.

\begin{definition}[Transversales Symmetriebrechungs-Ma{\ss} $S_{BC}$]
F\"ur den \"Ubergangswinkel $\phi\in[0,\pi/2]$ zwischen den Rotationszust\"anden
$B$ ($\rho(7)=2$) und $C$ ($\rho(11)=3$) ist das transversale Entropie-Proxy definiert durch
\[
S_{BC}(\phi)=\sin^2\phi\cdot\Lambda,
\quad\text{mit}\quad
\Lambda=\log(\mathrm{ratio})^2+\frac{\langle t_{BC}\rangle}{\langle t_{EA}\rangle}.
\]
\end{definition}

\begin{theorem}[Invariante TUR-Produkt-Schranke auf Schicht $[A]$]
\label{thm:tur-bc}
F\"ur den im $BC$-Unterraum generierten stochastischen Strom $J_{BC}$ mit mittlerer
Inkrement-Pr\"azision
$\mathcal{Q}_{BC}=\sigma_{J_{BC}}^2/\langle J_{BC}\rangle^2$
gilt unter dem Nullmodell $Q_0$:
\begin{equation}
\label{eq:tur-bc}
\mathcal{Q}_{BC}(\phi,\gamma)\cdot S_{BC}(\phi)
=
\frac{\Lambda}{T\cdot\gamma}
\ge
\mathcal{C}_{Q_0},
\end{equation}
wobei $\gamma=\mathrm{rate}_{BC}/\mathrm{rate}_{EA}$ die Ratenasymmetrie und $T$ das
Beobachtungsfenster bezeichnet.
\end{theorem}

\begin{proof}
Im station\"aren Zustand $p_B^*=p_C^*=1/2$ lauten der Erwartungswert und die Varianz
des Nettostroms $J_{BC}(T)$:
\[
\langle J_{BC}\rangle=T\cdot\gamma\sin\phi,\qquad
\sigma_{J_{BC}}^2=T\cdot\gamma.
\]
Daraus folgt f\"ur das Unsch\"arfema{\ss} $\mathcal{Q}_{BC}$:
\[
\mathcal{Q}_{BC}
=
\frac{T\cdot\gamma}{(T\cdot\gamma\sin\phi)^2}
=
\frac{1}{T\cdot\gamma\sin^2\phi}.
\]
Substitution von $\sin^2\phi=S_{BC}(\phi)/\Lambda$ liefert unmittelbar das
produktinvariante Resultat:
\[
\mathcal{Q}_{BC}\cdot S_{BC}(\phi)
=
\Bigl(\frac{1}{T\cdot\gamma\sin^2\phi}\Bigr)\cdot\bigl(\sin^2\phi\cdot\Lambda\bigr)
=
\frac{\Lambda}{T\cdot\gamma}.
\]
F\"ur $\phi\to 0$ divergiert $\mathcal{Q}_{BC}\to\infty$, w\"ahrend $S_{BC}(0)=0$;
die analytische Fortsetzung des Produkts kollabiert exakt auf den Invariantenwert
$\Lambda/(T\cdot\gamma)$, welcher durch die Nullmodell-Schranke $\mathcal{C}_{Q_0}$
nach unten beschr\"ankt ist (Eintrag $E$-ENT-$003$).
\end{proof}

\begin{remark}[Governance / Claim-Grenze]
\label{rem:claim-e-ent-003}
Gleichung~\eqref{eq:tur-bc} ist eine \emph{modellinterne} Invariante
(Projekt \emph{Entropie}, Tag \texttt{v0.1-tur-bc}; Lean-Pfad nur via
\texttt{EntropicBarrier} in \texttt{eabc-renorm} PR\,\#119, nicht als
\texttt{EEnt003} auf \texttt{master}).
Explizit \emph{nicht} behauptet wird eine W\"arme\-bad-Identifikation, ein globales
Optimierungsprinzip $Q^*$, oder ein formaler Collatz-/Survivor-Br\"uckenbeweis.
\end{remark}

\section{Zusammenfassung}

Die Zuordnung $\Phi$ fasst die sechs m\"oglichen L\"uckenklassen modulo $12$ in vier
geometrische Rotationszust\"ande zusammen:
\begin{align*}
\Phi(0) &= 0 \quad (\text{Identit\"at / Stillstand}),\\
\Phi(2)=\Phi(4) &= 1 \quad (\text{Vorw\"arts-Rotation um }+1),\\
\Phi(6) &= 2 \quad (\text{Reflexion / Halbdrehung um }+2),\\
\Phi(8)=\Phi(10) &= 3 \quad (\text{R\"uckw\"arts-Rotation um }-1\equiv 3).
\end{align*}
Im Zusammenspiel mit dem Invarianten-Eintrag $E$-ENT-$003$ zeigt sich, dass die
diskrete Ziffernkodierung auf $C_4$ \"uber die kontinuierliche TUR-Schranke
$\mathcal{Q}_{BC}\cdot S_{BC}\ge\mathcal{C}_{Q_0}$ eine modellintern abgesicherte,
informations- und bilanzierungskonsistente Repr\"asentation auf Schicht $[A]$ besitzt
(Bemerkung~\ref{rem:claim-e-ent-003}).

\end{document}